\documentclass[11pt]{article}
\usepackage[a4paper,margin=26mm]{geometry}
\usepackage[T1]{fontenc}
\usepackage{lmodern}
\usepackage{amsmath,amssymb,amsthm}
\usepackage{microtype}
\usepackage{array}
\usepackage{xurl}
\usepackage{hyperref}
\hypersetup{colorlinks=true,linkcolor=blue,urlcolor=blue,citecolor=blue,
 pdftitle={A product inequality for Walsh coefficients and two conditional Fourier entropy-influence bounds},
 pdfauthor={Horace Dong}}
\newtheorem{theorem}{Theorem}
\newtheorem{lemma}{Lemma}
\newtheorem{proposition}{Proposition}
\newtheorem{corollary}{Corollary}
\theoremstyle{definition}
\newtheorem{definition}{Definition}
\theoremstyle{remark}
\newtheorem{remark}{Remark}
\newcommand{\repoversion}{1.26.0}
\setlength{\emergencystretch}{1.5em}
\newcommand{\Z}{\mathbb Z}
\newcommand{\R}{\mathbb R}
\renewcommand{\lg}{\log_2}

\title{A product inequality for Walsh coefficients\\ and two conditional Fourier
entropy--influence bounds}
\author{Horace Dong\\ \small Independent researcher\\
\small\texttt{maxwelldhx@gmail.com}\\
\small ORCID: \href{https://orcid.org/0009-0002-7554-7548}{0009-0002-7554-7548}}
\date{23 November 2026}
\begin{document}
\sloppy
\maketitle

\begin{abstract}
The Fourier entropy--influence (FEI) conjecture of Friedgut and Kalai asks for
an absolute constant $C$ such that $H(f)\le C\,I(f)$ for every Boolean function
$f:\{-1,1\}^n\to\{-1,1\}$, where $H(f)$ is the entropy of the distribution
$p(S)=\hat f(S)^2$ and $I(f)=\sum_S|S|\,p(S)$ is the total influence. It is
open. We prove two partial results. Let $K$ be the number of nonzero Fourier
coefficients of $f$ and $t$ the number of nonempty sets among them. (a) If
$\delta>0$ and $p(S)\ge\delta/K$ for every $S$ with $p(S)\ne0$, then
$\delta\,(H(f)+2t/K)\le4\,I(f)$; in particular $H(f)\le4I(f)/\delta$, and
$H(f)\le4M^2I(f)$ when the nonzero coefficients differ in absolute value by a
factor at most $M$. (b) If $L\ge1$ and $p(S)\le L\,p(T)$ for all nonempty sets
$S\subseteq T$ with nonzero coefficients, then $H(f)\le4L\,I(f)$; the mass at
the empty set is unconstrained. Potapov proved the case of equal nonzero
magnitudes of~(a), with the constant $4$. The main ingredient is the inequality
$\sum_S\lg(1/p(S))+2t\le4\sum_S|S|$ over the nonzero coefficients, valid for
every Boolean function, which comes from an integer determinant. The constants
grow as $\delta\to0$ or $L\to\infty$, so these results do not approach the
conjecture itself. All six results of the paper (Theorems 1--5 and Corollary~1), including (a), (b) and the
product inequality, are formally verified in Lean~4. The
proofs were found and formalized by AI agents working under the author's
direction; their role is described at the end of the paper.
\end{abstract}

\section{Introduction}\label{sec:intro}

For $f:\{-1,1\}^n\to\{-1,1\}$ (any $n\ge0$) and $S\subseteq[n]=\{1,\dots,n\}$
let $\chi_S(x)=\prod_{i\in S}x_i$ and
$\hat f(S)=2^{-n}\sum_x f(x)\chi_S(x)$, the expectation under the uniform
measure. By Parseval the numbers $p(S)=\hat f(S)^2$ sum to $1$. The
\emph{Fourier entropy} and the \emph{total influence} are
\[
H(f)=\sum_{S}p(S)\lg\frac1{p(S)},\qquad I(f)=\sum_S|S|\,p(S),
\]
with $0\lg(1/0)=0$. The total influence equals the expected number of
coordinates $i$ for which flipping $x_i$ changes $f(x)$
(\cite[Section~2.1]{owz}, \cite[Theorem~2.38]{odo}); this identity is standard
and is not part of the formal verification, which uses the spectral form of
$I$. Friedgut and Kalai~\cite{fk} made the following conjecture; we quote the
quantifiers from O'Donnell, Wright and Zhou~\cite[Section~1]{owz}.

\begin{quote}
\textbf{FEI conjecture.} There is a constant $C$ such that $H(f)\le C\,I(f)$
for every $n$ and every $f:\{-1,1\}^n\to\{-1,1\}$.
\end{quote}

\noindent O'Donnell, Wright and Zhou proved it for symmetric functions, for
functions with symmetry group $S_{n_1}\times\dots\times S_{n_d}$ with $d$
bounded, and for read-once decision trees; they note that the AND and OR
functions force $C\ge4$, and that an explicit family forces $C\ge60/13$
(larger lower bounds are known today, e.g.\ $C>6.45$~\cite{hod}), and
they show the general bound $H(f)\le(1+\log n)I(f)+1$~\cite[Section~1]{owz}.
The conjecture is open in general; this paper does not prove it.

\paragraph{Notation.} Fix $f$. Let $U=\{S:\hat f(S)\ne0\}$ be its Fourier
support, $K=|U|\ge1$, $t=|U\setminus\{\varnothing\}|$ and
$W=\sum_{S\in U}|S|$. For $V\subseteq U$ put $W_V=\sum_{S\in V}|S|$ and
$t_V=|V\setminus\{\varnothing\}|$. A set $V\subseteq U$ is \emph{upward closed
in $U$} if $S\in V$, $T\in U$ and $S\subseteq T$ imply $T\in V$. The empty
family is upward closed. The empty set $\varnothing$ is the constant character,
$\hat f(\varnothing)=\mathbb E f$, and it belongs to $U$ unless $f$ is balanced.

\paragraph{The results.} All six statements below hold for every $n\ge0$ and
every $f$, without any balance or monotonicity assumption, and all six are
proved in Lean (Section~\ref{sec:lean}).

\begin{theorem}[signed product]\label{thm:prod}
$2^{2W}\prod_{S\in U}\hat f(S)=2^t z$ for a nonzero integer $z$.
\end{theorem}

\begin{theorem}[support entropy]\label{thm:supp}
$H(f)\le\lg K$, \ $K\lg K\le\sum_{S\in U}\lg\frac1{p(S)}$, \ and
\[
\sum_{S\in U}\lg\frac1{p(S)}+2t\ \le\ 4W .
\]
\end{theorem}

\begin{theorem}[lower mass bound]\label{thm:lower}
Let $\delta>0$ and suppose $\delta\le K\,p(S)$ for every $S\in U$. Then
\[
\delta\Bigl(H(f)+\frac{2t}{K}\Bigr)\ \le\ 4\,I(f).
\]
In particular $H(f)\le4I(f)/\delta$.
\end{theorem}

\noindent The hypothesis forces $\delta\le1$, by summing over $U$. Restricting
to $0<\delta\le1$ in the statement would change nothing; the Lean statement
requires only $\delta>0$.

\begin{corollary}[bounded ratio]\label{cor:ratio}
Let $M\ge1$ and suppose $|\hat f(S)|\le M|\hat f(T)|$ for all $S,T\in U$. Then
$H(f)\le4M^2I(f)$.
\end{corollary}

\begin{theorem}[upper sets]\label{thm:upper}
For every $V\subseteq U$ that is upward closed in $U$,
$2^{2W_V}\prod_{S\in V}\hat f(S)=2^{t_V}z$ for a nonzero integer $z$.
\end{theorem}

\begin{theorem}[comparable masses]\label{thm:comp}
Let $L\ge1$ and suppose $p(S)\le L\,p(T)$ for all nonempty $S,T\in U$ with
$S\subseteq T$. Then $H(f)\le4L\,I(f)$.
\end{theorem}

\noindent In Theorem~\ref{thm:comp} the hypothesis involves only nonempty sets:
the mass $p(\varnothing)$ is arbitrary, and sets that are not comparable under
inclusion are not compared. Taking $S=T$ shows that $L\ge1$ is automatic as
soon as $f$ is not constant. Theorem~\ref{thm:prod} is the case $V=U$ of
Theorem~\ref{thm:upper}; in the formal proof it is derived that way.

\paragraph{Scope and significance.} These are conditional bounds, and their
significance is modest. The constants $4/\delta$ and $4L$ grow without bound as
$\delta\to0$ or $L\to\infty$, and for many standard functions $\delta$ is
exponentially small: for the AND of $n\ge2$ bits, $K\min_U p=4\cdot2^{-n}$. A
product of $m$ AND functions on disjoint blocks of three bits has
$K\min_Up=2^{-m}$ and a ratio $9^{m-1}$ between two nonempty comparable masses,
while its entropy and influence both grow only linearly in $m$ (an observation
in the Codex notes, not formalized). So neither condition holds with a uniform
parameter for all Boolean functions, and the two theorems give no information
on the general conjecture. What they give is a dimension-free bound on two
classes of non-flat spectra, and the product inequality of
Theorem~\ref{thm:supp}, which holds for all $f$.

\paragraph{Prior work and what is new.} Potapov~\cite{pot} proved $H(f)\le4I(f)$
for all plateaued functions, that is, functions whose nonzero Fourier
coefficients all have the same absolute value~\cite[Theorem~2]{pot}, and $2$
for partially bent functions~\cite[Theorem~1]{pot}; he states that $4$ is sharp
for plateaued functions. For a plateaued $f$, $K\,p(S)=1$ on $U$, so
Theorem~\ref{thm:lower} with $\delta=1$ and Corollary~\ref{cor:ratio} with
$M=1$ give $H(f)\le4I(f)$ again: this case is due to Potapov, and our
Theorem~\ref{thm:lower} contains no new information on it beyond the additive
term $2t/K$. His proof also uses the matrix of the inclusion order on the
Fourier support, which is unitriangular, but over the field with two elements:
he deduces from $\det(MM^{\mathsf T})=1$ a permutation $\pi$ of the support
such that the number of support points above both $y$ and $\pi(y)$ is odd for
every $y$~\cite[Lemma~1]{pot}, and combines it with a bound on
the algebraic degree of the support indicator of a plateaued
function~\cite[Proposition~4 and Theorem~3]{pot}. We put the actual signed
coefficients on the diagonal and use integrality of the determinant over
$\Z$, so that the magnitudes may differ. We do not claim the triangular matrix
as a new idea. Other recent work on FEI includes Arunachalam, Chakraborty,
Koucký, Saurabh and de Wolf~\cite{acksw} (min-entropy, parity certificates and
flat approximating polynomials), Gonz\'alez, MacManus and Pereyra~\cite{gmp}
(stopping binary trees and tribes-like classes) and Shmalo~\cite{shm} (linear
threshold functions, compositions and selectors), and Paraashar and
Saurabh~\cite{ps} (approximate spectral norms and parity certificates); none
of these considers a
lower bound on the nonzero masses, a comparability condition on them, or a
product of all nonzero coefficients. The closest comparison for
Theorem~\ref{thm:comp} that the searches found is Shmalo's manuscript on
equivalent forms of the conjecture~\cite{shm2}, which bounds $H(f)$ by the
expected logarithm of the ratio between the mass of all supersets of $S$ and
$p(S)$; our hypothesis bounds this ratio only from below, and that manuscript
does not state a bound under it. What is new here, as far as the searches
described at the end found, is Theorem~\ref{thm:supp} (with
Theorems~\ref{thm:prod} and~\ref{thm:upper}), Theorem~\ref{thm:lower} for
$\delta<1$, Corollary~\ref{cor:ratio} for $M>1$, and Theorem~\ref{thm:comp}.

\section{Finite differences and the signed product}\label{sec:prod}

\begin{lemma}[finite differences]\label{lem:tail}
For every $R\subseteq[n]$ the number $2^{|R|}\sum_{T\supseteq R}\hat f(T)$ is an
integer, and it is even if $R\ne\varnothing$.
\end{lemma}

\begin{proof}
For $A\subseteq R$ let $x^{(A)}$ be the point with $x_i=-1$ exactly for
$i\in A$. By Fourier inversion,
\[
\sum_{A\subseteq R}(-1)^{|A|}f(x^{(A)})=\sum_T\hat f(T)\sum_{A\subseteq R}(-1)^{|A|}(-1)^{|A\cap T|}
=\sum_T\hat f(T)\prod_{i\in R}\bigl(1-(-1)^{\mathbf 1[i\in T]}\bigr),
\]
and the product is $2^{|R|}$ if $R\subseteq T$ and $0$ otherwise. So
$2^{|R|}\sum_{T\supseteq R}\hat f(T)$ is a signed sum of $2^{|R|}$ values $\pm1$:
an integer, and even when $R\ne\varnothing$.
\end{proof}

\begin{proof}[Proof of Theorem~\ref{thm:upper} (and of Theorem~\ref{thm:prod})]
Let $V$ be upward closed in $U$ and order $V$ by nondecreasing cardinality.
The matrix $Z_{S,T}=\mathbf 1[S\subseteq T]$ ($S,T\in V$) is upper triangular
with ones on the diagonal, so $\det Z=1$. Let $D=\operatorname{diag}(2^{|S|})$,
$C=\operatorname{diag}(\hat f(S))$ and $B=DZCZ^{\mathsf T}D$. Then
\[
B_{S,R}=2^{|S|+|R|}\sum_{T\in V,\ T\supseteq S\cup R}\hat f(T).
\]
A set $T\in U$ with $T\supseteq S\cup R$ contains $S\in V$, so it lies in $V$;
hence the sum runs over all $T\supseteq S\cup R$ (the others have
$\hat f(T)=0$). Since $|S|+|R|\ge|S\cup R|$, Lemma~\ref{lem:tail} shows that
$B_{S,R}$ is an integer, and an even one whenever $S\ne\varnothing$. Dividing
the $t_V$ rows with $S\ne\varnothing$ by $2$ gives an integer matrix $B'$, and
\[
\det B'=2^{-t_V}\det B=2^{-t_V}\,2^{2W_V}\prod_{S\in V}\hat f(S).
\]
It is an integer, and nonzero because every $\hat f(S)$ with $S\in V$ is
nonzero. For $V=\varnothing$ both sides are $1$. Theorem~\ref{thm:prod} is the
case $V=U$.
\end{proof}

\section{Entropy consequences}\label{sec:entropy}

\begin{lemma}[logarithms]\label{lem:log}
If $V\subseteq U$ satisfies the conclusion of Theorem~\ref{thm:upper}, then
$\sum_{S\in V}\lg\frac1{p(S)}+2t_V\le4W_V$.
\end{lemma}

\begin{proof}
Take $\lg$ of absolute values: $2W_V+\sum_V\lg|\hat f(S)|=t_V+\lg|z|\ge t_V$,
and $\lg(1/p(S))=-2\lg|\hat f(S)|$.
\end{proof}

\begin{proof}[Proof of Theorem~\ref{thm:supp}]
By concavity of $\lg$ with weights $p(S)$,
$H(f)\le\lg\sum_Up(S)\cdot\frac1{p(S)}=\lg K$. By convexity of $-\lg$ with
uniform weights on $U$ and $\frac1K\sum_Up(S)=\frac1K$,
$\frac1K\sum_U\lg\frac1{p(S)}\ge\lg K$. The third inequality is
Lemma~\ref{lem:log} with $V=U$ and Theorem~\ref{thm:prod}.
\end{proof}

\begin{proof}[Proof of Theorem~\ref{thm:lower}]
From $\delta\le Kp(S)$ we get $\delta W\le K\sum_U|S|p(S)=K\,I(f)$. By
Theorem~\ref{thm:supp} and $\delta>0$,
\[
K\delta\Bigl(H(f)+\frac{2t}K\Bigr)=\delta\bigl(KH(f)+2t\bigr)\le\delta(K\lg K+2t)
\le\delta\Bigl(\sum_U\lg\tfrac1{p(S)}+2t\Bigr)\le4\delta W\le4K\,I(f).
\]
Divide by $K$. Since $2t/K\ge0$, also $\delta H(f)\le4I(f)$.
\end{proof}

\begin{proof}[Proof of Corollary~\ref{cor:ratio}]
Let $m=\min_Up$. Then $p(S)\le M^2m$ on $U$, so $1=\sum_Up\le KM^2m$, and
$\delta=M^{-2}$ satisfies the hypothesis of Theorem~\ref{thm:lower}.
\end{proof}

\section{Comparable masses}\label{sec:comp}

\begin{lemma}[layers]\label{lem:layer}
Let $A$ be a finite set, $a,g:A\to\R$ with $a\ge0$, and suppose
$\sum_{i\in A,\ a_i\ge r}g_i\le0$ for every $r>0$. Then $\sum_{i\in A}a_ig_i\le0$.
\end{lemma}

\begin{proof}
Let $0<r_1<\dots<r_k$ be the positive values of $a$ and $r_0=0$. Then
$\sum_ia_ig_i=\sum_{j=1}^k(r_j-r_{j-1})\sum_{a_i\ge r_j}g_i\le0$.
\end{proof}

\begin{lemma}[the empty frequency]\label{lem:dc}
For $0\le p\le1$, \ $p\lg\frac1p\le2(1-p)$.
\end{lemma}

\begin{proof}
For $p=0$ both sides are $\ge0$ with the left side $0$. For $0<p\le1$,
$\ln(1/p)\le1/p-1$ gives $p\ln(1/p)\le1-p$, and $\ln2\ge\frac12$ (from
$\ln x\ge1-1/x$ at $x=2$) gives $p\lg(1/p)\le(1-p)/\ln2\le2(1-p)$.
\end{proof}

\begin{proof}[Proof of Theorem~\ref{thm:comp}]
Let $U'=U\setminus\{\varnothing\}$. If $U'=\varnothing$, then
$p(\varnothing)=1$ and $H(f)=I(f)=0$. Otherwise define, for $S\in U'$,
\[
a_S=\max\{p(T):T\in U',\ T\subseteq S\}.
\]
Then $p(S)\le a_S\le L\,p(S)$ by the hypothesis, and $a$ is nondecreasing under
inclusion on $U'$. For $r>0$ the set $V_r=\{S\in U':a_S\ge r\}$ is upward
closed in $U$ (a set of $U$ containing a member of $V_r$ is nonempty and has a
larger value of $a$), and $t_{V_r}=|V_r|$. By Theorem~\ref{thm:upper} and
Lemma~\ref{lem:log}, $\sum_{S\in V_r}g_S\le0$ with
$g_S=\lg\frac1{p(S)}+2-4|S|$. Lemma~\ref{lem:layer} gives
\[
\sum_{S\in U'}a_S\Bigl(\lg\frac1{p(S)}+2\Bigr)\le4\sum_{S\in U'}a_S|S|.
\]
On the left $\lg(1/p(S))+2\ge0$, so $a_S$ may be decreased to $p(S)$; on the
right $a_S$ may be increased to $Lp(S)$. This gives
$\sum_{U'}p(S)\lg\frac1{p(S)}+2(1-p(\varnothing))\le4L\,I(f)$ by Parseval, and
Lemma~\ref{lem:dc} applied to $p(\varnothing)$ (which is $0$ if
$\varnothing\notin U$) completes the proof.
\end{proof}

\section{Examples and remarks (not formalized)}\label{sec:examples}

Nothing in this section is part of the Lean verification.

\begin{remark}[sharpness of $4$ in Theorem~\ref{thm:supp}]\label{rem:address}
The address function on $r$ address bits and $2^r$ data bits returns the data
bit selected by the address. Its Fourier support consists of the sets formed by
a subset of the address bits and one data bit, all with $p(S)=4^{-r}$; so
$K=4^r$, $\lg K=2r$ and the mean of $|S|$ over $U$ is $r/2+1$
(cf.~\cite[Sections~5.7 and~8]{shm}; in Section~5.7 $H=2r$ and $I=1+r/2$ are stated). Hence the
ratio $\lg K\big/\bigl(\frac1K\sum_U|S|\bigr)$ tends to $4$, and the constant $4$
in the consequence $K\lg K\le4W$ of Theorem~\ref{thm:supp} cannot be lowered.
The same family shows that $4$ cannot be lowered in Theorem~\ref{thm:lower} at
$\delta=1$ (Potapov states this sharpness for plateaued
functions~\cite[Section~2]{pot}, using a construction of Li and Tang; the
address function gives another such family). Our own
computation reproduced these values for $r=1,2,3$.
\end{remark}

\begin{remark}[the constant in Theorem~\ref{thm:comp} for $L=1$]\label{rem:and}
For the AND function of $n$ bits, all nonempty coefficients have the same
absolute value $2^{1-n}$, so Theorem~\ref{thm:comp} applies with $L=1$; a direct
computation gives $H/I\to4$ as $n\to\infty$ (we found $3.84$ for $n=8$ and $3.91$
for $n=12$). So the constant $4$ in Theorem~\ref{thm:comp} cannot be lowered at
$L=1$. Since the best constant in the FEI conjecture is
larger than $4$ ($60/13$ in~\cite[Section~1]{owz}, more than $6.45$
in~\cite{hod}), the hypothesis of Theorem~\ref{thm:comp} with
$L=1$ cannot hold for all Boolean functions. The computation for the AND
function was added while writing this paper and was not part of the documents
read by the reviewers; the observation that a lower bound above $4$ for the FEI
constant rules out $L=1$ for all functions was made by the Codex reviewer of
the extension.
\end{remark}

\begin{remark}[the two conditions are independent]\label{rem:selectors}
Codex's written proofs give selector examples, built on the selector identities
of~\cite[Section~8]{shm}. For $r\ge3$ and $N=2^r$, split the $N$ addresses into
one cell of $N/2-1$ addresses and $N/2+1$ singleton cells, give each cell its
own data bit, and output the data bit of the cell containing the address. All
cells have odd size, every magnitude $N^{-1}|\sum_{z\in C}\chi_S(z)|$ is
nonzero, $K\min_Up=\frac12+\frac2N>\frac12$, and the ratio of the largest to
the smallest nonzero magnitude is $N/2-1\to\infty$. So
Theorem~\ref{thm:lower} gives $H\le8I$ on this family while
Corollary~\ref{cor:ratio} gives no uniform bound. If the large cell is chosen as
the hyperplane $z_1=0$ without the zero address and the output is multiplied by
the parity of all address bits, the masses become nondecreasing under inclusion
on the support ($L=1$), and products of independent copies keep $L=1$ while
$K\min_Up$ tends to $0$; the untwisted examples have $K\min_Up>\frac12$ but
$L\to\infty$ (with $L=(N/2-1)^2$ for the cells used in our computation). So
neither condition implies the other (the product
observation is due to a Codex reviewer). Our computation for $r=3,4$
reproduced $K=48,160$, $K\min_Up=0.75,0.625$, magnitude ratios $3,7$, and
$L=9,49$ for the untwisted and $L=1$ for the twisted examples. We do not claim
that these selector examples are new.
\end{remark}

\section{Formal verification}\label{sec:lean}

{\sloppy
The six statements are formally verified in Lean~4 (version 4.34.1) with
Mathlib~\cite{mathlib} (commit \texttt{d13f23b7}). The eight source files
\texttt{Challenge}, \texttt{UpperChallenge}, \texttt{Walsh},
\texttt{Determinant}, \texttt{Analysis}, \texttt{Layer},
\texttt{UpperAnalysis} and \texttt{Check} are in version~\repoversion{} of the
repository~\cite{repo}, directory \texttt{lean/Research/TaskV/FEI/}; records of
the freezing and of the checks are in the folder
\texttt{certificates/fei-nonflat-spectra}.

\paragraph{The statements.} The statement file \texttt{Challenge.lean}, in ASCII
transliteration (the \texttt{import Mathlib} line, a two-line docstring, blank
lines, and
the lines \texttt{namespace Fei}, \texttt{open Finset} and \texttt{end Fei} are
omitted; the file uses Unicode: \texttt{Nat}, \texttt{Int}, \texttt{Real} stand
for $\mathbb N$, $\Z$, $\R$, \texttt{delta} for $\delta$, \texttt{sum} and
\texttt{prod} for the big operators $\sum$ and $\prod$, \texttt{in},
\texttt{!=}, \texttt{<=}, \texttt{->}, \texttt{/\char`\\}, \texttt{forall},
\texttt{exists} and \texttt{subseteq} for $\in$, $\ne$, $\le$, $\to$, $\wedge$,
$\forall$, $\exists$ and $\subseteq$):
{\scriptsize
\begin{verbatim}
abbrev Cube (n : Nat) := Finset (Fin n)
noncomputable def value {n : Nat} (f : Cube n -> Bool) (x : Cube n) : Real :=
  if f x then -1 else 1
noncomputable def character {n : Nat} (s x : Cube n) : Real :=
  prod i in x, if i in s then (-1 : Real) else 1
noncomputable def coeff {n : Nat} (f : Cube n -> Bool) (s : Cube n) : Real :=
  (sum x : Cube n, value f x * character s x) / 2 ^ n
noncomputable def support {n : Nat} (f : Cube n -> Bool) : Finset (Cube n) := by
  classical exact univ.filter fun s => coeff f s != 0
noncomputable def mass {n : Nat} (f : Cube n -> Bool) (s : Cube n) := (coeff f s) ^ 2
noncomputable def entropy {n : Nat} (f : Cube n -> Bool) : Real :=
  sum s in support f, mass f s * (-Real.logb 2 (mass f s))
noncomputable def influence {n : Nat} (f : Cube n -> Bool) : Real :=
  sum s in support f, mass f s * s.card
noncomputable def weightSum {n : Nat} (f : Cube n -> Bool) : Nat :=
  sum s in support f, s.card
noncomputable def nonemptyCount {n : Nat} (f : Cube n -> Bool) : Nat :=
  ((support f).erase {}).card
def SignedProduct : Prop :=
  forall (n : Nat) (f : Cube n -> Bool), exists z : Int, z != 0 /\
    (2 : Real) ^ (2 * weightSum f) * (prod s in support f, coeff f s) =
      (2 : Real) ^ nonemptyCount f * (z : Real)
def SupportEntropy : Prop :=
  forall (n : Nat) (f : Cube n -> Bool),
    entropy f <= Real.logb 2 (support f).card /\
    (support f).card * Real.logb 2 (support f).card <=
      sum s in support f, -Real.logb 2 (mass f s) /\
    (sum s in support f, -Real.logb 2 (mass f s)) + 2 * nonemptyCount f <=
      (4 : Real) * weightSum f
def LowerMassFEI : Prop :=
  forall (n : Nat) (f : Cube n -> Bool) (delta : Real), 0 < delta ->
    (forall s in support f, delta <= (support f).card * mass f s) ->
    delta * (entropy f + (2 : Real) * (nonemptyCount f : Real) /
      ((support f).card : Real)) <= 4 * influence f
def BoundedRatioFEI : Prop :=
  forall (n : Nat) (f : Cube n -> Bool) (M : Real), 1 <= M ->
    (forall s in support f, forall t in support f, |coeff f s| <= M * |coeff f t|) ->
    entropy f <= 4 * M ^ 2 * influence f
\end{verbatim}}
\noindent The line break inside \texttt{LowerMassFEI} is ours, and \texttt{\{\}}
stands for the empty set $\varnothing$. A point of \texttt{Cube n} is the set of
its coordinates equal to $-1$, and the value \texttt{true} means $-1$; this is a
bijection with the functions $\{-1,1\}^n\to\{-1,1\}$, for every $n$ including
$n=0$. \texttt{character s x} is $(-1)^{|s\cap x|}=\chi_s(x)$, \texttt{coeff}
is $\hat f(s)$, \texttt{support f} is $U$ (the exact set of nonzero
coefficients), \texttt{mass} is $p$, and \texttt{Real.logb 2} is $\lg$. Summing
\texttt{entropy} over the support only is the convention $0\lg(1/0)=0$, and
\texttt{influence} is the spectral form of $I(f)$. \texttt{weightSum} is $W$ and
\texttt{nonemptyCount} is $t$. So \texttt{SignedProduct},
\texttt{SupportEntropy}, \texttt{LowerMassFEI} and \texttt{BoundedRatioFEI}
are Theorems~\ref{thm:prod}, \ref{thm:supp}, \ref{thm:lower} and
Corollary~\ref{cor:ratio}; the hypothesis of \texttt{LowerMassFEI} is written
$\delta\le K\,p(S)$ without division, and its conclusion is the stronger form
with $2t/K$. The second statement file \texttt{UpperChallenge.lean} (imports
\texttt{Fei.Challenge}, in the repository
\texttt{Research.TaskV.FEI.Challenge}; same omissions) contains exactly
{\footnotesize
\begin{verbatim}
def UpperProduct : Prop :=
  forall (n : Nat) (f : Cube n -> Bool) (V : Finset (Cube n)),
    V subseteq support f ->
    (forall s in V, forall w in support f, s subseteq w -> w in V) ->
    exists z : Int, z != 0 /\
      (2 : Real) ^ (2 * sum s in V, s.card) * (prod s in V, coeff f s) =
        (2 : Real) ^ (V.erase {}).card * (z : Real)
def ComparableMassFEI : Prop :=
  forall (n : Nat) (f : Cube n -> Bool) (L : Real), 1 <= L ->
    (forall s in (support f).erase {}, forall t in (support f).erase {}, s subseteq t ->
      mass f s <= L * mass f t) ->
    entropy f <= 4 * L * influence f
\end{verbatim}}
\noindent These are Theorems~\ref{thm:upper} and~\ref{thm:comp}; in
\texttt{ComparableMassFEI} the empty set is removed from both quantified
indices, so $p(\varnothing)$ is unconstrained. None of the statements encodes
the general FEI conjecture, the identity between $I$ and the flip influence, or
the examples of Section~\ref{sec:examples}.

\paragraph{Review and freezing.} Both statement files were compiled with no
diagnostics before they were frozen: \texttt{Challenge.lean} between 03:58:30
and 03:58:38~UTC on 21 November 2026 and \texttt{UpperChallenge.lean} between
04:02:10 and 04:02:15~UTC. Each was compared with the written proof by a Codex
agent that had not written the statement or the proof (approvals recorded at
03:58:15 and 04:00:58~UTC); the first reviewer had flagged an ambiguous
natural-number expression in an earlier draft of \texttt{LowerMassFEI}, which
was replaced by the real quotient shown above. The files were frozen by
recording their SHA-256 hashes, \texttt{b277a582\ldots} at 03:58:46~UTC and
\texttt{0731c4d4\ldots} at 04:02:30~UTC, and were not changed afterwards. The
development receipts record no compilation of a proof file before the first
freeze (the first is at 04:01:22~UTC); the extension statement was written at
04:00:40~UTC and approved at 04:00:58~UTC; \texttt{Walsh} and
\texttt{Analysis} (which contains the lemma \texttt{upper\_log} used for
Theorem~\ref{thm:comp}) were last changed at 04:01:38 and 04:01:54~UTC, after
that approval but before the second freeze, and \texttt{Layer},
\texttt{UpperAnalysis} and the final version of \texttt{Determinant} were
created or rewritten after the second freeze. \texttt{Analysis},
\texttt{Layer}, \texttt{UpperAnalysis} and \texttt{Determinant} were first
compiled successfully at 05:03--05:05~UTC.

\paragraph{The proofs.} The check file \texttt{Check.lean} does not have the
form \texttt{theorem check\_X : Challenge.X := Proof.X} used in our other
papers: each check theorem has the frozen proposition as its type and composes
proof lemmas in its body. Apart from the \texttt{import} lines and the lines
\texttt{namespace Fei}, \texttt{open Finset} and \texttt{end Fei}, it contains
exactly
{\footnotesize
\begin{verbatim}
theorem check_upperProduct : UpperProduct := upperProduct_of (fun _ => tail_integral)
theorem check_signedProduct : SignedProduct := signedProduct_of_upper check_upperProduct
theorem check_supportEntropy : SupportEntropy :=
  supportEntropy_of check_signedProduct (fun _ => parseval)
theorem check_lowerMassFEI : LowerMassFEI :=
  lowerMassFEI_of check_supportEntropy (fun _ => parseval)
theorem check_boundedRatioFEI : BoundedRatioFEI :=
  boundedRatioFEI_of check_lowerMassFEI (fun _ => parseval)
theorem check_comparableMassFEI : ComparableMassFEI :=
  comparableMassFEI_core (fun _ => upper_layer_sum) (fun _ => parseval)
    (fun n f V hs hu => upper_log f V hs (check_upperProduct n f V hs hu))
\end{verbatim}}
\noindent together with three \texttt{example} lines that evaluate
\texttt{character} and \texttt{coeff} on tiny cases and six \texttt{\#print
axioms} commands. The checked declarations are
\texttt{Fei.check\_signedProduct}, \texttt{Fei.check\_supportEntropy},
\texttt{Fei.check\_lowerMassFEI}, \texttt{Fei.check\_boundedRatioFEI},
\texttt{Fei.check\_upperProduct} and \texttt{Fei.check\_comparableMassFEI}.
Every module is in the namespace \texttt{Fei}, so a redeclaration of a frozen
name would be a Lean error; no proof module declares one, and the final
reviewer confirmed with its own file that
\texttt{example : Fei.X := Fei.check\_X} compiles for all six. Table~\ref{tab:lean}
lists the Lean counterparts of the steps (file names in parentheses, without
\texttt{.lean}; names marked $^*$ are \texttt{private}).

\begin{table}[!htbp]
\centering\footnotesize
\renewcommand{\arraystretch}{1.05}
\begin{tabular}{@{}>{\raggedright\arraybackslash}p{0.34\textwidth}>{\raggedright\arraybackslash}p{0.62\textwidth}@{}}
\hline
Step of the written proof & Lean (file) \\
\hline
Orthogonality, inversion, Parseval & \texttt{orthogonal}$^*$, \texttt{inversion}$^*$,
\texttt{parseval} (\texttt{Walsh}) \\
Lemma~\ref{lem:tail} & \texttt{signed\_face}$^*$, \texttt{face\_sum}$^*$,
\texttt{tail\_integral} (\texttt{Walsh}) \\
Theorem~\ref{thm:upper}: $\det Z=1$, integer entries, row factors &
\texttt{inc\_det}$^*$, \texttt{entry\_factor}$^*$, \texttt{row\_prod}$^*$,
\texttt{upperProduct\_of} (\texttt{Determinant}) \\
Theorem~\ref{thm:prod} from Theorem~\ref{thm:upper} &
\texttt{signedProduct\_of\_upper} (\texttt{Determinant}) \\
Lemma~\ref{lem:log} & \texttt{upper\_log} (\texttt{Analysis}) \\
Theorem~\ref{thm:supp} (two Jensen steps) & \texttt{entropy\_bounds}$^*$,
\texttt{supportEntropy\_of} (\texttt{Analysis}) \\
Theorem~\ref{thm:lower}, Corollary~\ref{cor:ratio} & \texttt{lowerMassFEI\_of},
\texttt{boundedRatioFEI\_of} (\texttt{Analysis}) \\
Lemma~\ref{lem:layer} & \texttt{upper\_layer\_sum} (\texttt{Layer}) \\
Lemma~\ref{lem:dc}, Theorem~\ref{thm:comp} & \texttt{entropy\_dc\_le}$^*$,
\texttt{comparableMassFEI\_core} (\texttt{UpperAnalysis}) \\
\hline
\end{tabular}
\caption{Lean counterparts of the steps of the proof.}
\label{tab:lean}
\end{table}

\noindent The Jensen steps are proved from $\ln x\ge1-1/x$ and $\ln x\le x-1$,
and the majorant $a_S$ of Section~\ref{sec:comp} is defined in
\texttt{comparableMassFEI\_core} as a maximum over the nonempty supported
subsets of $S$ together with $S$ itself. The axioms of all six check theorems
are exactly \texttt{propext}, \texttt{Classical.choice} and
\texttt{Quot.sound}. No file uses \texttt{sorry}, \texttt{native\_\allowbreak
decide}, \texttt{bv\_\allowbreak decide}, \texttt{decide} or an \texttt{axiom}
declaration; the only option set is one \texttt{maxHeartbeats 800000} in
\texttt{UpperAnalysis}.

\paragraph{Checks.} The nine modules (the eight files and a root module
\texttt{Fei} that imports \texttt{Check}) were compiled module by module in two
new empty directories on the author's Mac, one after the other, with the frozen
Lake configuration and warnings treated as errors, on 21 November 2026
(05:05:26--05:11:43~UTC); each directory also ran one build of the whole
package, replayed every module with \texttt{leanchecker} and printed the axioms.
All commands ended with exit code~$0$, the \texttt{leanchecker} runs printed
nothing, the two directories produced identical \texttt{.olean} files, and the
largest memory use of a single step was $6.0$~GB. An attempt started ten seconds
earlier stopped at the version probe, before building any module, because the
timing tool was not allowed to read a system clock setting; it was kept on
record. The final reviewer (see the next section) compiled the nine modules
again with plain \texttt{lean} in a new directory of its own, together with a
short audit file of its own, and replayed every module with
\texttt{leanchecker} (21 November 2026, between 17:03 and 17:09~UTC); all ten
modules compiled with exit code~$0$ and no warnings, and all
\texttt{leanchecker} runs ended with exit code~$0$ and empty output. Mathlib and
its dependencies were taken from a shared prebuilt cache at the pinned commits
and were not rechecked; the \texttt{leanchecker} replays cover the project
modules. In the repository the eight files have the module names
\texttt{Research.TaskV.FEI.Challenge}, \dots,
\texttt{Research.TaskV.FEI.Check}, keep the namespace \texttt{Fei} (so the
checked names \texttt{Fei.check\_X} are unchanged), and differ from the checked
files only in their \texttt{import} lines; \texttt{Challenge.lean}, which
imports only Mathlib, is byte-identical (SHA-256 \texttt{b277a582\ldots}). The
root module \texttt{Fei} of the checked package, which only imports
\texttt{Check}, is not part of the repository.
These versions were compiled and replayed with \texttt{leanchecker} again, module by module, in two new empty directories on the author's Mac (22 November 2026, 21:36--21:39~UTC), with no warnings and identical \texttt{.olean} files and the same axiom lists.
In the whole-project rebuild of version~\repoversion{} of the repository (on the author's Windows computer, in one pass from 22 to 23 November 2026), these files were compiled and replayed with \texttt{leanchecker} again together with the rest of the project; no module needed a second attempt (details in the repository). The formalization has not been checked with comparator or with an
independent kernel.

\par}

\section*{How this paper was produced}
This work was carried out from 21 to 23 November 2026 (UTC) in a workflow
directed by the author, using OpenAI Codex with the model GPT-6.2-Astra
(reasoning level ultra) and Anthropic Claude Code with the model Claude
Opus~5.5. The author asked for a search for well-known open problems, and
Claude Code passed this request to Codex at 03:19~UTC on 21 November; in that
search Codex prepared a note on the FEI conjecture among others. The author
then asked that Codex continue to attack it, and Claude Code started a Codex
session for this at 03:53~UTC; its task allowed partial results with a complete
written proof approved by an agent that had not derived it to go on to a frozen
statement and a Lean proof. According to the file times, by 03:56~UTC a Codex
agent had written the proof of Theorems~\ref{thm:prod}--\ref{thm:lower} and
Corollary~\ref{cor:ratio}, and another Codex agent that had not taken part in
it approved it at 03:57~UTC; the extension to upper sets
(Theorems~\ref{thm:upper} and~\ref{thm:comp}) was written by 03:58~UTC and
approved by a non-author Codex reviewer at 04:00~UTC, who also checked all
$278$ functions with $n\le3$ and $3{,}460$ upper sets with its own program
(review and freezing: Section~\ref{sec:lean}). The Codex session reached its usage limit at 04:04~UTC;
Claude Code resumed it at 05:03~UTC, and the formalization and the checks were
completed by 05:12~UTC. The handoff notes record that a reviewer of the
written proofs later took part in the Lean work on the extension, so the final
review was assigned to an agent that had written none of the files. The Lean proofs
follow the written proofs.

The final review was done by a Claude Code agent (model Claude Opus~5.5) that
had written none of the files, on 21 November 2026 from 17:03~UTC. It read the
source~\cite{owz} itself, compared the statement files with it, read the written
proofs and all proof files, checked the hashes and the receipts, rebuilt the
formalization as described above, and tested the six statements with its own
exact program on all $278$ functions with $n\le3$ and $3{,}998$ random functions
with $4\le n\le7$, finding no violation. Its decision was APPROVE\_FORMAL, with
no blocking or major remarks; its six minor remarks concern the written account
and the form of the check file and are taken into account here (among them the
form of \texttt{Check.lean}, the marking of material that is not in Lean, the
stronger form of Theorem~\ref{thm:lower}, the unconstrained empty frequency in
Theorem~\ref{thm:comp}, and the search before publication). Codex then wrote an
English summary of the mathematics (17:36~UTC). Claude Code (Claude Opus~5.5)
checked the sources, recomputed in Python Theorems~\ref{thm:prod}, \ref{thm:lower}
and~\ref{thm:comp} on $1{,}178$ functions
(all functions with $n\le3$ and random ones with $n=4,5,6$), the examples of
Remarks~\ref{rem:address}--\ref{rem:selectors} and the AND function, and wrote
this text from Codex's manuscripts and the Lean files; the AND computation in
Remark~\ref{rem:and} and the scope paragraph of Section~\ref{sec:intro} were added while writing.
An independent Claude Code agent (model Claude Opus~5.5) that had written none of the files reviewed this paper, and two others reviewed the release; their corrections were applied. Throughout, ``we'' refers to this workflow; the reading and
searching described in the paper were done by the agents. Codex agents searched
for prior work on 21 November (03:53--04:03~UTC: arXiv, GitHub, MathDB, the
Moonshot and Epoch problem lists, repositories of AI-assisted mathematics and
general web searches, with a separate search for the upper-set condition); some
citation indexes refused queries (rate limit). The final reviewer searched
again on 21 November, and on 22 November 2026 (21:11~UTC) Claude Code queried
the arXiv for new preprints on the conjecture and found none after~\cite{ps} (4 November
2026), with~\cite{pot} and~\cite{shm} still in their first versions. Up to
22~November 2026 these searches found no earlier statement of the inequality of
Theorem~\ref{thm:supp}, of the product statements of Theorems~\ref{thm:prod}
and~\ref{thm:upper}, of Theorem~\ref{thm:lower} for $\delta<1$, of
Corollary~\ref{cor:ratio} for $M>1$, or of Theorem~\ref{thm:comp} (details in
the repository). On 23~November 2026 (11:10--11:28~UTC) Claude Code searched again (the arXiv, the OpenAlex and Semantic Scholar lists of works citing the main sources, Zenodo, DataCite, GitHub and repositories of AI-assisted mathematics); the only new related work found is two notes of Y.~Shmalo posted on 22 and 23~November 2026, on symmetric functions and on functions of bounded degree; their entropy--influence bounds assume symmetry or a bound on the degree (functions of degree at most~$d$ satisfy the hypothesis of Theorem~\ref{thm:lower} with $\delta=4^{1-d}$, and on them the second note's constant, of order $2\lg d$, is far smaller than the constant $4^d$ given by Theorem~\ref{thm:lower}), and neither note contains the results here. This reports the coverage of the searches,
not a guarantee that they were complete. The author is not a professional
mathematician. The author has read this paper and, using a side-by-side table, compared its Lean statements with Theorems~\ref{thm:prod}--\ref{thm:comp} and Corollary~\ref{cor:ratio}. No human mathematician has read the proofs or
reviewed this paper. The correctness of the six statements rests on the Lean
proofs and on the agreement of the Lean statements with the definitions
of~\cite{owz}; the written proofs are given for the reader. AI systems are not
authors of this paper; the author takes full responsibility for its content.

\medskip
\noindent Thanks are due to Vladimir N. Potapov for the plateaued case and the
triangular matrix on which this work builds, to Ryan O'Donnell, John Wright and
Yuan Zhou for the framework, and to Yitzchak Shmalo for the selector identities.

{\footnotesize
\begin{thebibliography}{99}
\setlength{\itemsep}{0pt}
\bibitem{acksw} S.~Arunachalam, S.~Chakraborty, M.~Kouck\'y, N.~Saurabh and
R.~de~Wolf, Improved bounds on Fourier entropy and min-entropy, \emph{ACM
Trans.\ Comput.\ Theory} 13 (2021), no.~4; arXiv:1809.09819v2 (2021).
\bibitem{repo} H.~Dong, \emph{ai4math-results}, repository,
\url{https://github.com/Horace-Maxwell/ai4math-results}, version~\repoversion{}
(2026); archived at Zenodo under the concept DOI
\url{https://doi.org/10.5281/zenodo.22970254}.
\bibitem{fk} E.~Friedgut and G.~Kalai, Every monotone graph property has a sharp
threshold, \emph{Proc.\ Amer.\ Math.\ Soc.} 124 (1996), no.~10, 2993--3002,
\url{https://doi.org/10.1090/S0002-9939-96-03732-X}.
\bibitem{gmp} M.~J. Gonz\'alez, P.~MacManus and M.~C. Pereyra, Further evidence
towards the Fourier entropy-influence conjecture, arXiv:2606.00246v2 (2026).
\bibitem{hod} R.~Hod, Improved lower bounds for the Fourier entropy/influence
conjecture via lexicographic functions, arXiv:1711.00762v1 (2017).
\bibitem{mathlib}
The mathlib Community, The Lean mathematical library, in \emph{Proc.\ 9th ACM
SIGPLAN Int.\ Conf.\ Certified Programs and Proofs} (CPP '20), ACM, 2020,
367--381, \url{https://doi.org/10.1145/3372885.3373824}.
\bibitem{odo} R.~O'Donnell, \emph{Analysis of Boolean Functions}, Cambridge
University Press, 2014, \url{https://doi.org/10.1017/CBO9781139814782}.
\bibitem{owz} R.~O'Donnell, J.~Wright and Y.~Zhou, The Fourier entropy--influence
conjecture for certain classes of Boolean functions, in \emph{Automata,
Languages and Programming} (ICALP 2011, Part~I), Lecture Notes in Comput.\ Sci.\
6755, Springer, 2011, 330--341, \url{https://doi.org/10.1007/978-3-642-22006-7_28}.
\bibitem{ps} M.~Paraashar and N.~Saurabh, Fourier entropy bounds via
approximations, arXiv:2611.06451v1 (2026).
\bibitem{pot} V.~N. Potapov, On the Fourier entropy-influence conjecture for
Boolean plateaued functions, arXiv:2609.21563v1 (2026).
\bibitem{shm} Y.~Shmalo, Fourier entropy of threshold functions and their
compositions, arXiv:2610.16069v1 (2026).
\bibitem{shm2} Y.~Shmalo, Four equivalent forms of the Fourier entropy-influence
conjecture, manuscript,
\url{https://github.com/yspennstate/fei-equivalent-forms} (commit
\texttt{305c4f8e}, 2026).
\end{thebibliography}}
\end{document}
